
% ============================================================================
\clearpage
\section*{v4.10.0-r62 Source-Only: Octonion Copies in the Sedenions and the Remainder Generation (k321--k330, H138--H147)}
\addcontentsline{toc}{section}{v4.10.0-r62 Source-Only: Octonion Copies in the Sedenions and the Remainder Generation (k321--k330, H138--H147)}
\label{sec:v4100-r62-sedenion-copies}
% ============================================================================

\paragraph{Purpose.}
On 25 November 2026 the author asked whether the sedenions explain why only three or four copies arise for the generations, and conjectured that
beyond them only arrangements of zero divisors remain.  Ten tests followed (H138--H147, scripts k321--k330), each fixed in an advance note before
the computation and recorded as a research note with scripts and outputs.  This layer collects them, each with its status as recorded at the time,
on the author's decision of 26 November 2026 to bring the Companion up to date.  It adds no new computation; every statement below is checked in
the cited scripts.  Where a part of a test was not blind, or was found after the fact, this is said.  The tests H128--H137 (k278--k320) of the same
days belong to other lines and are not collected here.  Earlier layers are not rewritten.

\noindent\textbf{Source-only status.}
This v4.10.0-r62 layer is source-only.  It does not build a PDF, does not assign a new DOI,
and does not authorize a public release.  The v5.4.0 Version DOI
\href{https://doi.org/10.5281/zenodo.24500009}{10.5281/zenodo.24500009} and the Concept DOI
\href{https://doi.org/10.5281/zenodo.19210728}{10.5281/zenodo.19210728} are unchanged by it.

\subsection*{v4.10.0-r62 setting}
\label{subsec:v4100-r62-setting}

The tower is the Cayley--Dickson tower of r59 with \(HS(3)=\mathbb O\), \(HS(4)=\mathbb S\) (sedenions, \(A_4\)), \(HS(5)=A_5\), \(HS(6)=A_6\).
\(G_2=\operatorname{Aut}(\mathbb O)\) acts blockwise on \(A_n=\mathbb O\otimes\mathbb R^{2^{n-3}}\).  A \emph{copy} is a unital, \(G_2\)-stable,
eight-dimensional subalgebra \(O_v=\mathbb R\,1\oplus\operatorname{Im}\mathbb O\otimes v\) with a unit direction \(v\); it is \emph{octonionic} if
it is alternative.  This definition is a setting; colour-stable subalgebras (only \(su(3)\), H130) are more numerous and are not counted.
Colour is the stabiliser \(su(3)\subset g_2\) of \(e_7\); the charge of a copy is the k85 charge \(Q_v\) built from \(L_u,R_u\) with
\(u=e_7\otimes v\), with spectrum \(0,\tfrac13^3,\tfrac23^3,1\) on \(C_v=\mathbb C\otimes O_v\).  Because \(u\) depends on the sign of \(v\), the
\emph{orientation} of every copy enters; in r62.5--r62.8 it is transported by the automorphisms that permute the copies.  The global sign of the
orientation is not set here: by the author's rule it is to follow from the field equation.  Reading copies as generations is a reading that is
tested below, not assumed.  Tools: Python (numpy, scipy) for searches and all selections, Wolfram Language exactly for \(n=4,5\) and for one
representative of each class; Qiskit was not used (no circuit and no measurement sequence occurs).

\subsection*{v4.10.0-r62 results}
\label{subsec:v4100-r62-results}

\paragraph{Theoremlet v4.10.0-r62.1 (three copies per doubling; H138, k321; \(HS(4)\) not blind, \(HS(5)\), \(HS(6)\) blind).}
For \(v=(t_1,t_2)\) the closure condition of \(O_v\) in \(\mathbb S\) is \(t_2(3t_1^2-t_2^2)=0\), that is
\[
\operatorname{Im}\bigl((t_1+i t_2)^3\bigr)=0 .
\]
It has exactly three solution lines (\(0^\circ,60^\circ,120^\circ\)); a fourth \(G_2\)-stable octonion copy does not exist in \(\mathbb S\)
(Wolfram exact).  The threefold count comes from the conjugation in the doubling product.  The three copies span only two directions and no two
are orthogonal: three labels, two independent carriers, the same linear dependence as recorded for K28 (rank 6 instead of 9).  Higher levels:
\begin{center}
\small
\begin{tabular}{@{}lllll@{}}
\toprule
level & closed \(O_v\) & octonionic & rank of directions & largest orthogonal set (all / octonionic)\\
\midrule
\(HS(4)\) & 3 (exact) & 3 & 2 & 1 / 1\\
\(HS(5)\) & 15 (exact) & \(9=3\times3\) & 4 & 2 / 1\\
\(HS(6)\) & \(\ge 207\) (search not saturated) & \(\ge 79\) & 8 & \(\ge5\) / 1\\
\bottomrule
\end{tabular}
\end{center}
On \(HS(5)\) the nine octonionic copies are exactly the products \(v=a\otimes b\) of the three sedenion directions.  H138a (\(N_5=9\),
\(N_6=27\)) and H138b (at most three orthogonal copies) are \emph{refuted} in the counting fixed in advance.  Orthogonal, hence independent,
copies are never octonionic on any level computed.  \emph{Status:} ``three, not four'' is shown within one doubling over \(\mathbb O\), conditional
on the definition of a copy; that the three copies are three generations is not shown, and the number of generations remains an input.

\paragraph{Theoremlet v4.10.0-r62.2 (zero divisors between and beyond the copies; H139, k322; partly not blind).}
For \(x=c\,e_1+s\,e_{10}\),
\[
\det L_x=(c-s)^4(c+s)^4(c^2+s^2)^4 ,
\]
zero exactly for \(c^2=s^2\); for the parallel pair \(x=c\,e_1+s\,e_9\) it is \((c^2+s^2)^8\) and never zero (Wolfram exact).  In the copy
language an element of \(\mathbb S\) is a \(7\times2\) block \([a\ b]\), and the zero divisors are exactly the elements with \(a\perp b\) and
\(|a|=|b|\): equal weight \(\tfrac12\) in every copy direction (84 basis diagonals and 500 samples, deviation \(\le6\cdot10^{-15}\); H139b, H139c
confirmed, the latter not blind).  An element of a single direction is never a zero divisor.  On \(HS(5)\) the six closed, non-alternative
algebras beyond the nine octonionic copies contain no zero divisor among themselves, so H139a is \emph{refuted by its wording}; found after the
fact (k322c, k321c): every purely imaginary element of them is a zero divisor of all of \(A_5\) (kernel of dimension 12), and the same holds on
\(HS(6)\) within rounding.  The nine octonionic copies contain no zero divisor.  The author's conjecture holds in this form: what goes beyond the
\(3\times3\) octonionic copies are families of zero divisors, leaking outward, not inward.  It does not follow that zero divisors represent or
forbid a fourth generation.

\paragraph{Theoremlet v4.10.0-r62.3 (same colour, own charge, no common charge; H140, k323; partly not blind).}
The three copies in \(\mathbb S\) have the same \(su(3)\) (principal-angle cosine \(1\) to twelve digits), which commutes with the \(S_3\) of the
copies, and each carries one generation: \(Q_k\) maps \(C_k\) into itself with spectrum \(0,\tfrac13^3,\tfrac23^3,1\), and the cyclic automorphism
\(\psi\) sends \(Q_k\) to \(Q_{k+1}\) (H140a, H140b confirmed).  A charge that preserves all three copies does not exist in the family
\(\alpha+\beta L_x+\gamma R_x\): the leak is exactly
\[
\bigl(12|\beta|^2+12|\gamma|^2-18\operatorname{Re}\bar\beta\gamma\bigr)\bigl(|x|^2+x_8^2\bigr)\ \ge\ 3|x|^2
\]
(Wolfram exact; H140c confirmed).  The three copies together cover 15 complex dimensions instead of 24 (H140d, not blind).  By the consequence
fixed in advance the copies in \(\mathbb C\otimes\mathbb S\) are therefore \emph{not} three generations; they are three equivalent readings of one.

\paragraph{Theoremlet v4.10.0-r62.4 (the nine missing states; H141, k324; partly not blind).}
Take the copies as separate carriers \(V=C_0\oplus C_{120}\oplus C_{240}\) (24 states) with \(\Phi(x_0,x_1,x_2)=x_0+x_1+x_2\) into
\(\mathbb C\otimes\mathbb S\).  The image has dimension 15, the cokernel is \(e_8\), and the kernel \(K\) (states that cancel in \(\mathbb S\)) has
dimension 9 with \(su(3)\) content \(1\oplus1\oplus1\oplus3\oplus\bar3\).  The compressed charge on \(K\) has spectrum
\(\tfrac13^3,\tfrac12^3,\tfrac23^3\) (Wolfram exact).  The colour part \(3\oplus\bar3\) is charge stable and carries the quark charges; the three
singlets carry \(\tfrac12\) and leak (leak \(\tfrac12\)).  \(K\) is not closed under the product, its components are not zero divisors, and
\([Q_\oplus,P_K]\ne0\): no common charge (H141a--d confirmed).  A fourth member built from the remainder has a stable quark part but no leptons.

\paragraph{Theoremlet v4.10.0-r62.5 (four and five copies on \(HS(5)\): the remainder generation; H142, H143, k325, k326; partly not blind).}
Among the 126 selections of four of the nine octonionic copies on \(HS(5)\), with transported orientation, exactly 9 have a kernel that
contains a part with spectrum \(0,\tfrac13^3,\tfrac23^3,1\), exactly charge stable (leak \(\le4\cdot10^{-15}\)) and decomposing as
\(1\oplus1\oplus3\oplus\bar3\): a complete \emph{remainder generation}, leptons included, plus two leaking singlets with \(\tfrac12\).  H142b, which
predicted no such selection, is \emph{refuted}.  The nine selections are those whose directions satisfy an alternating relation
\(v_1-v_2+v_3-v_4=0\); none contains all three sedenion copies.  A first run failed its control because the directions carried arbitrary signs;
the orientation rule was added to the advance note before the evaluated run, and this is recorded.  With five copies (H143) the rule
\[
n_{\rm gen}=\dim R^0,\qquad R^0=\Bigl\{c:\ \textstyle\sum_k c_kv_k=0,\ \sum_k c_k=0\Bigr\},
\]
holds in 2016 of 2016 cases (126 selections \(\times\) 16 orientations).  At most one remainder generation occurs (45 selections with transported
orientation), and each stems from exactly one of the nine relations of H142.  The fifth copy adds only a further quark set or leaking states with
compressed charges \(\tfrac12\), \(\tfrac1{10}\), \(\tfrac9{10}\), \((5\mp\sqrt{10})/10\), \((10\mp\sqrt{10})/20\).  No common charge exists
(minimal squared leak \(3.5\) to \(6.0\)); no kernel component is a zero divisor.  Because the orientation enters the relations, the count depends
on the open sign; the maximum 1 does not.  Whether remainder states are physical states (they cancel in \(A_5\)) is not shown.

\paragraph{Theoremlet v4.10.0-r62.6 (the remainder carries no sector of its own; H144, k327; parts a, b not blind).}
Setting: the factor \(\mathbb C\) of \(\mathbb C\otimes A_5\) is the C-factor of the frames as in r51, with sector quantity
\(N=\sum_a x_a^2\) or \(s=\sum_{a\ge1}x_a^2\).  Under either quantity the remainder generation has the signature of a copy generation: quarks
\(s=0\), leptons \(|s|=\tfrac12\) (Wolfram exact).  H144c, which predicted \(|s|<10^{-6}\) for the remainder leptons, is \emph{refuted}.  The
criterion of H144d could not fail (recognised after the fact).  The frame maps act within each carrier, while every remainder state lies on at
least three carriers, so they do not connect copies and remainder.  Side result: time reversal \(z\mapsto\bar z\) maps the remainder generation
to itself with \(q\mapsto1-q\), which is the global sign change of the orientation (\(1-Q_v=Q_{-v}\)).  The open global sign is therefore the
choice of the time direction (particle against antiparticle); it is left to the field equation.

\paragraph{Theoremlet v4.10.0-r62.7 (the remainder is built on \(HS(0)\) but is no algebra; H145, k328; largely not blind).}
With \(D^0=\{c:\sum_kc_k=0\}\) and \(R=\{c:\sum_kc_kv_k=0\}\), in all 3360 selection--orientation cases with 3, 4, 5 copies
\[
K=(D^0\otimes e_0)\oplus(R\otimes\operatorname{Im}\mathbb O),\qquad \dim K=(m-1)+7\dim R
\]
(projector distance \(\le2\cdot10^{-14}\); Wolfram exact for one selection).  With transported orientation the remainder generation is
\(R^0\otimes\mathbb C\mathbb O\): a whole octonion copy over a real coefficient line \(c=(-1,-1,1,1)\), on which charge and colour do not act.
The line has the marks of \(HS(0)\): real, only \(\pm c\) as a choice, fixed under time reversal.  The remainder is not an \(HS(0)\) as an algebra:
it is not closed (for the H142 selections the product of two remainder states lies entirely outside), and \(\sum_kc_k^3=0\).  H145c is
\emph{refuted by its wording} (it holds for transported orientation only); with mixed signs a copy carries the mirrored octonion, and H145a for the
stable part needed a corrected form found after the fact.  The formula explains the rule of r62.5 analytically.  No new tower starts from the line.

\paragraph{Theoremlet v4.10.0-r62.8 (no masses from the copy structure; H146, k329; parts a, b not blind).}
Setting: the Gram readout of r14--r16, \(\sqrt{m_k}\) the eigenvalues of the Gram matrix of the generation states,
\(K=\sum m/(\sum\sqrt m)^2\).  Three copies give at most \(K=19/32\) (leptons) and \(\tfrac12\) (colour) over all 84 selections; the sedenion
triple gives \(\sqrt m=(\tfrac34,\tfrac34,\tfrac32)\), \(K=\tfrac38\).  With four copies in an H142 relation every charge sector has exactly one
eigenvalue \(0\), with eigenvector exactly the line \(c\) of r62.7; the other three give \(K=59/128\) (leptons) and \(11/32\) (colour), with
\(|K-\tfrac23|\ge0.206\) and mass ratios at most \(11.1\) against the measured \(3477\) (H146a--c confirmed).  For charged leptons the frame mask
is a similarity transformation (form S, r57), so the verdict does not depend on the frame.  The copy structure alone gives no masses.  In this
readout the remainder generation has eigenvalue zero; this agrees with the author's statement that on \(HS(0)\) and \(HS(1)\) there is no mass in
the strict sense, and is not an independent derivation of it.

\paragraph{Theoremlet v4.10.0-r62.9 (the \(HS(0)\) line is a scalar, not a spin-2 graviton; H147, k330; largely not blind).}
The line \(c\) is the exact null vector of the Gram matrix in all 54 generation selections and is neutral.  Under \(SO(3)\subset SU(3)_c\) (all
eight totally real coordinate triples) \(\mathbb C\mathbb O=2\times\text{spin }0\oplus2\times\text{spin }1\), and the line itself has spin \(0\);
under the principal subgroup of \(G_2\) (representation theory only) \(1\oplus7\) is spin \(0\oplus\) spin \(3\).  Spin 2 occurs in the outer
product \(\operatorname{Sym}^2(\mathbb O)\), where the symmetric algebra product vanishes on it (norm \(\le7\cdot10^{-15}\)), and in the derivation
channel \(D_{x,y}\), which reaches the spin-2 part of \(su(3)\) fully (rank 5) but is antisymmetric in the two states (H147a--d confirmed; one
evaluation of part c was corrected and the correction is recorded).  The line is thus a neutral scalar, dilaton-like, not the graviton.  A
symmetric spin-2 coupling would need the metric as an additional rule; coupling strength, Newton's constant and fifth-force bounds are not
computed.  \emph{Exploration, after the fact:} the remainder has two scalar directions \(c\otimes e_0\) and \(c\otimes e_7\); the charge rotation
\(\exp(i\varphi Q)\) turns \(e_0\) into \(e_7\) by \(\varphi/2\), so the family \(\cos\theta\,e_0+\sin\theta\,e_7\) has a \emph{free} mixing angle
\(\theta\) under the charge symmetry, and no automorphism mixes the two.  Whether the \(e_7\) end is a Higgs is not shown.  This angle is the
starting point of r63.

\subsection*{v4.10.0-r62 frames}
\label{subsec:v4100-r62-frames}

\begin{center}
\small
\begin{tabular}{@{}>{\raggedright\arraybackslash}p{0.34\textwidth}>{\raggedright\arraybackslash}p{0.22\textwidth}>{\raggedright\arraybackslash}p{0.36\textwidth}@{}}
\toprule
item & frame & comparison\\
\midrule
r62.1--r62.5, r62.7 (algebra, charge) & M (field value) & none across frames\\
r62.6 (sector quantity) & M of the copies; T by \(z\mapsto iz\); time reversal \(z\mapsto\bar z\) & \(s\mapsto-s\) resp.\ \(s\mapsto\bar s\); \(|s|\) unchanged, verdict frame independent\\
r62.8 (Gram readout) & M & form S leaves eigenvalues and \(K\) unchanged\\
r62.9 (spin content) & M; choice of the spatial \(SO(3)\) is a setting & none\\
P & not used & ---\\
\bottomrule
\end{tabular}
\end{center}
The tachyon map \(z\mapsto iz\) and time reversal \(z\mapsto\bar z\) are kept apart throughout (r51, r81).

\subsection*{v4.10.0-r62 register entries}
\label{subsec:v4100-r62-register}

\begin{center}
\small
\begin{tabular}{@{}>{\raggedright\arraybackslash}p{0.40\textwidth}>{\raggedright\arraybackslash}p{0.52\textwidth}@{}}
\toprule
statement & status after v4.10.0-r62\\
\midrule
three octonion copies per doubling (\(\operatorname{Im}z^3=0\)) & shown, conditional on the definition of a copy\\
H138a, H138b; H139a; H142b; H144c; H145c & refuted (H139a, H145c by wording)\\
H139b, c; H140a--d; H141a--d; H142a, c, d; H143a--d; H144a, b; H145a (for \(K\)), b, d; H146a--c; H147a--d & confirmed, with the blindness recorded per part\\
copies \(=\) three generations in \(\mathbb C\otimes\mathbb S\) & not the case (no common charge, 15 of 24 states)\\
number of generations & remains an input\\
remainder generation as physical state & not shown\\
orientation and global sign & open; to follow from the field equation, not set\\
\(HS(0)\) line as graviton & no (scalar); scalar gravity not computed\\
T1 & remains a conjecture\\
OP1, OP3, OP4, OP5 & open\\
\bottomrule
\end{tabular}
\end{center}

\subsection*{v4.10.0-r62 notes}
\label{subsec:v4100-r62-notes}

{\raggedright
\texttt{notes/\allowbreak research/\allowbreak hilbert\_\allowbreak shift\_\allowbreak felder\_\allowbreak 2026\_\allowbreak 11\_\allowbreak 24/\allowbreak generationszahl\_\allowbreak h138/}:
advance notes \texttt{Vorab\_\allowbreak Notiz\_\allowbreak H138\_\allowbreak *.md} to \texttt{Vorab\_\allowbreak Notiz\_\allowbreak H147\_\allowbreak *.md}, result notes
\texttt{Ergebnis\_\allowbreak H138\_\allowbreak *.md} to \texttt{Ergebnis\_\allowbreak H147\_\allowbreak *.md}, scripts \texttt{k321} to \texttt{k330} (each with a
Wolfram cross-check \texttt{b}, some with \texttt{c}, \texttt{d}) and their actual outputs \texttt{k3NN\_\allowbreak out.txt}.
\par}

\subsection*{v4.10.0-r62 ledger and guard update}
\label{subsec:v4100-r62-ledger-guard}

\[
R_{\rm chk}^{v4100\text{-}r62}=(1,1,0)
\]
for (status recorded; computations checked in the cited scripts; no new computation in this layer; generation number, orientation and masses not derived).

The layer spends no guard:
\[
R_{\rm glob}^{(v4100\text{-}r62,+)}=(0,1,1,1,1),
\qquad
\Gamma_{\rm guard}^{v4100\text{-}r62}=(1,1,1,1,1),
\qquad
\beta_4=1/256\ \text{unspent}.
\]
The layer does not imply exact finite \(C5\) zero closure, exact all-readout terminality, \(C6/T2\), a global photon-clock/RH theorem,
or global all-readout closure.  OP1--OP5, C5 and C6/T2 remain open.
